\cgeentry{cge_microsim}{Pass a simulation to household survey data}

\cgesec{Description}

\cmd{cge_microsim} passes the results of a simulation to the household survey
in memory, top-down and without feedback to the model. Each survey household
belongs to a household group of the model and keeps its own income
composition and, if they are given, its own budget shares. Each income
component changes as the model says that its source changes for the
household's group; nominal consumption follows the household's income and its
group's direct tax and savings rates; and consumption is deflated by a price
index at the prices households pay, including VAT. The command reports real
consumption per person, the poverty headcount and poverty gap, and the Gini
coefficient, before and after the simulation, for all households and for
each group, and can add the household's new consumption and price index to
the data.

\cgesec{Quick start}

Pass simulation \cmd{notariff} to the survey in memory, in which \cmd{grp}
holds each household's model group and \cmd{cons} its consumption; every
household's consumption changes as its group's does
\begin{stcode}
cge_microsim, sim(notariff) group(grp) consumption(cons)
\end{stcode}
As above, with sampling weights \cmd{w}, consumption per person and a poverty
line of 1,200 per person
\begin{stcode}
cge_microsim [pweight = w], sim(notariff) group(grp) consumption(cons) size(hhsize) line(1200)
\end{stcode}
Households' own incomes from wages, profits, remittances and grants, and
their own budget shares for food, goods and services
\begin{stcode}
cge_microsim [pweight = w], sim(notariff) group(grp) consumption(cons) size(hhsize) ///
    income(wage:lab profit:cap remit:row grants:gov)                              ///
    budget(food:com_agr goods:com_man services:com_ser) line(1200)
\end{stcode}
As above, and add the variables \cmd{nt_nom}, \cmd{nt_price}, \cmd{nt_real}
and \cmd{nt_pc}, replacing them if they exist
\begin{stcode}
cge_microsim [pweight = w], sim(notariff) group(grp) consumption(cons) size(hhsize) ///
    income(wage:lab profit:cap remit:row grants:gov) line(1200) generate(nt) replace
\end{stcode}

\cgesec{Syntax}

\begin{stsyntax}
\cmdb{cge_microsim} [\argu{if}] [\argu{in}] [\argu{weight}]\texttt{,}
\optarg{}{sim}{name} \optarg{}{group}{varname} \optarg{cons}{umption}{varname}
[\argu{options}]
\end{stsyntax}

\begin{opttable}
\optgroup{Main}
\optarg{}{sim}{name} & the simulation; required\\
\optarg{}{group}{varname} & each household's model group; required\\
\optarg{cons}{umption}{varname} & household consumption; required\\
\optarg{}{size}{varname} & household size; default 1\\
\optarg{inc}{ome}{var\texttt{:}source \ldots} & income components and their sources\\
\optarg{bud}{get}{var\texttt{:}commodity \ldots} & budget items and the commodities whose prices apply\\
\optarg{}{line}{\#} & poverty line per person\\
\optgroup{Results}
\optarg{gen}{erate}{prefix} & add new consumption and the price index as variables\\
\opt{}{replace} & allow the variables of \cmd{generate()} to be replaced\\
\opt{nod}{isplay} & suppress the table of results\\
\end{opttable}

\cmd{pweight}s, \cmd{aweight}s, \cmd{fweight}s and \cmd{iweight}s are
allowed; all are used as proportional household weights.

\cgesec{Options}

\begin{optdesc}
\item[\optarg{}{sim}{name}] names the simulation whose results are passed to
the survey: a simulation solved by \manref{cge_solve}, a period of a run of
\manref{cge_dynamic}, or a solution imported by \mbox{\manref{cge_gams}}. Changes
are measured from the calibrated base.

\item[\optarg{}{group}{varname}] gives each household's group: a string
variable whose values are labels of the model's households, or a numeric
variable whose value labels, or failing those whose values, are. A value that
is not a household of the model is an error.

\item[\optarg{cons}{umption}{varname}] is the household's total consumption,
the welfare measure, in the survey's units.

\item[\optarg{}{size}{varname}] is the household's size, which must be
positive. Real consumption, poverty and inequality are then measured per
person, over persons. The default is a size of 1.

\item[\optarg{inc}{ome}{var\texttt{:}source \ldots}] lists the household's
income components, numeric variables, each with its source:
\begin{itemize}\setlength{\itemsep}{0pt}\setlength{\labelsep}{0.6em}
\item a factor of the model: the group's income from that factor, or the
factor's income if the group has none in the base;
\item \cmd{gov}: transfers from the government, set in real terms and paid at the consumer price index;
\item \cmd{row}: transfers from abroad, such as remittances;
\item \cmd{dom}: transfers from other domestic institutions, such as
dividends;
\item \cmd{cpi}: income indexed to the consumer price index;
\item \cmd{fixed}: nominal income that does not change.
\end{itemize}
Without \cmd{income()}, every household's nominal consumption changes as its
group's does.

\item[\optarg{bud}{get}{var\texttt{:}commodity \ldots}] lists spending on
budget items, numeric and not negative, and the commodity whose price applies
to each. An item may instead name an activity, for home consumption of its
output at the activity's price. A household's price index then uses its own
budget shares; without \cmd{budget()}, or for a household whose items sum to
zero, it uses its group's base budget shares, including home consumption.

\item[\optarg{}{line}{\#}] is the poverty line per person, in the units of
consumption. Without it, or with a line that is not positive, no poverty
measures are computed.

\item[\optarg{gen}{erate}{prefix}] adds the variables \argu{prefix}\cmd{_nom}
(new nominal consumption), \argu{prefix}\cmd{_price} (the household's price
index, relative to the base), \argu{prefix}\cmd{_real} (real consumption, at base
prices) and \argu{prefix}\cmd{_pc} (real consumption per person). They are
missing for observations outside the sample. \argu{prefix} may have at most
26 characters.

\item[\opt{}{replace}] allows the variables of \cmd{generate()} to be
replaced. Without it, a variable that exists is an error.

\item[\opt{nod}{isplay}] suppresses the table of results.
\end{optdesc}

\cgesec{Remarks and examples}

A model has a few household groups; poverty and inequality are measured on
household surveys. \mbox{\cmd{cge_microsim}} links the two in the top-down way
described by Bourguignon, Bussolo and Pereira da Silva (2008): the
model gives, for each group, the change in each source of income and in each
price; the survey gives each household's income composition, budget shares,
size and weight. The survey's own income and consumption data drive the
results, so a household that lives on wages gains more than its group when
wages rise relative to other income. There is no feedback from the survey to
the model. Every substantial source of income should be mapped, and the
survey and the SAM should refer to the same year where possible.

An observation is used if it satisfies \argu{if} and \argu{in} and has no
missing value in the group, consumption, weight, size or the variables of
\cmd{income()} and \cmd{budget()}.

\cgesubsec{Example 1: every household follows its group}

The synthetic SAM is calibrated as in \manref{cge_solve}, and every import
tariff is removed, with the direct tax rates of households adjusting to keep
government savings fixed:
\stlog{ms_sim}
The survey is synthetic, generated with a fixed seed: 1,000 households of one
to six persons in the model's single group, \cmd{hhd}, with weights, incomes
from wages, profits and government transfers, and spending on food, goods
and services:
\stlog{ms_survey}
\cgeneed{16\baselineskip}
Without \cmd{income()} and \cmd{budget()}, each household's nominal
consumption changes as its group's, and its price index is the group's:
\stlog[\scriptsize]{ms_group}
The table shows, for all households and for each group in the survey, mean
real consumption per person, weighted by persons, in the base and in the
simulation, and its per cent change; the poverty headcount and poverty gap,
in per cent; and the Gini coefficient of consumption per person, in per cent.
Every household of a group changes in the same proportion, so with one group
the Gini coefficient cannot change; the poverty measures change because
households near the line cross it.

\cgesubsec{Example 2: households' own incomes and budgets}

With \cmd{income()}, each household's income components change as their
sources do for its group; with \cmd{budget()}, each household has its own
price index:
\stlog[\scriptsize]{ms_income}
Real consumption per person falls by less, 0.389 against 0.431 per cent, and
the Gini coefficient now changes, from 36.32 to 36.33 per cent, because
households with different shares of wages, profits and transfers, and of
food, goods and services, are affected differently. The survey and the SAM
are synthetic: the results illustrate the command and are not findings about
any country.

\cgesubsec{Example 3: stored results and new variables}

The measures for all households are returned as scalars, and the whole table,
with the severity of poverty (FGT2), in \res{r(table)}; poverty measures and
Gini coefficients are stored as proportions, not per cent:
\stlog{ms_return}
\cmd{generate(nt)} added each household's new nominal consumption, price
index, real consumption and real consumption per person:
\stlog{ms_vars}
These variables can be analysed with any Stata command, for example to
compute other poverty or inequality measures, or measures by region.

\cgesec{Stored results}

\cmd{cge_microsim} stores the following in \cmd{r()}:
\begin{storedres}
\resgroup{Scalars}
\res{r(N)} & number of households used\\
\res{r(fgt0_base)} & poverty headcount in the base, all households (missing without \cmd{line()})\\
\res{r(fgt0_sim)} & poverty headcount in the simulation, all households\\
\res{r(gini_base)} & Gini coefficient in the base, all households\\
\res{r(gini_sim)} & Gini coefficient in the simulation, all households\\
\resgroup{Matrices}
\res{r(table)} & one row for all households (\cmd{all}) and one for each group in the survey; columns \cmd{mean_base}, \cmd{mean_sim}, \cmd{mean_pct}, \cmd{fgt0_base}, \cmd{fgt0_sim}, \cmd{fgt1_base}, \cmd{fgt1_sim}, \cmd{fgt2_base}, \cmd{fgt2_sim}, \cmd{gini_base} and \cmd{gini_sim}\\
\end{storedres}

\cgesec{Methods and formulas}

Household $i$ belongs to group $h = h(i)$ and has income components
$y_{ik}$ with sources $k$. The model gives, for each group, the ratio of the
simulated to the base value of each source,

\begin{cgetable}{p{0.20\linewidth}p{0.42\linewidth}p{0.30\linewidth}}
\toprule
Source & Ratio $\rho_{hk}$ & If the group's base value is zero\\
\midrule\endhead
\bottomrule\endlastfoot
factor $f$ & $\V{YIF}_{hf} / \V{YIF}^0_{hf}$ & $\V{YF}_f / \V{YF}^0_f$\\
\cmd{gov} & $\V{trnsfr}_{h,gov}\,\V{CPI} / (\V{trnsfr}^0_{h,gov}\,\V{CPI}^0)$ & $\V{CPI}/\V{CPI}^0$\\
\cmd{row} & $\V{trnsfr}_{h,row}\,\V{EXR} / (\V{trnsfr}^0_{h,row}\,\V{EXR}^0)$ & $\V{EXR}/\V{EXR}^0$\\
\cmd{dom} & $\sum_d \V{shii}_{hd}\,\V{DY}_d / \sum_d \V{shii}^0_{hd}\,\V{DY}^0_d$ & $\V{CPI}/\V{CPI}^0$\\
\cmd{cpi} & $\V{CPI}/\V{CPI}^0$ & \\
\cmd{fixed} & 1 & \\
\end{cgetable}
where $\V{DY}_d = (1 - \V{MPS}_d)(1 - \V{TINS}_d)\,\V{YI}_d$ is the income
that institution $d$ passes on or spends. The household's nominal
consumption becomes
\begin{equation*}
C^1_i = C^0_i\,\frac{\sum_k \rho_{hk}\,y_{ik}}{\sum_k y_{ik}}\,\kappa_h, \qquad
\kappa_h = \frac{(1 - \V{MPS}_h)(1 - \V{TINS}_h)}{(1 - \V{MPS}^0_h)(1 - \V{TINS}^0_h)},
\end{equation*}
and $C^1_i = C^0_i\,\V{EH}_h / \V{EH}^0_h$ for a household whose components
sum to zero, and for every household without \cmd{income()}. The household's
price index is
\begin{equation*}
P_i = \sum_j s_{ij}\,\frac{p_j}{p^0_j},
\end{equation*}
over commodities and home-consumed activity outputs $j$, where $p_j$ is the
price households pay for commodity $j$, $\V{PH}_j$, including VAT, or the
activity price $\V{PA}_j$ for home consumption; $s_{ij}$ are the household's
budget shares from \cmd{budget()} or else the group's base shares,
$p^0_j\,\V{QH}^0_{jh}/\V{EH}^0_h$ (with $\V{QHA}^0_{jh}$ for home
consumption). Real consumption is $R_i = C^1_i/P_i$, a first-order
(Laspeyres) approximation for households with LES demand, and with household
size $n_i$ and weight $w_i$, consumption per person is $r^0_i = C^0_i/n_i$ in
the base and $r^1_i = R_i/n_i$ in the simulation.

Over the households of a row of the table, with person weights
$v_i = w_i n_i$, the mean is $\sum_i v_i r_i / \sum_i v_i$, and for the
poverty line $z$
\begin{equation*}
\text{FGT}_\alpha = \frac{\sum_i v_i\,\mathbf{1}[r_i < z]\,\bigl((z - r_i)/z\bigr)^\alpha}{\sum_i v_i},
\qquad \alpha = 0, 1, 2 .
\end{equation*}
The Gini coefficient is computed from the Lorenz curve over persons: with
households sorted by $r_i$, $F_i$ the cumulative share of person weights and
$L_i$ the cumulative share of $v_i r_i$, and $F_0 = L_0 = 0$,
\begin{equation*}
G = 1 - \sum_i (F_i - F_{i-1})(L_i + L_{i-1}).
\end{equation*}
\Chapref{chap:microsim} describes the link in full.

\cgesec{References}

Bourguignon, F., M. Bussolo and L. A. Pereira da Silva, eds. 2008.
\textit{The Impact of Macroeconomic Policies on Poverty and Income
Distribution: Macro-Micro Evaluation Techniques and Tools}. World Bank and
Palgrave Macmillan.

\cgesec{Also see}

\begin{alsosee}
\seeentry{cge_solve}{Solve the model}
\seeentry{cge_report}{Standard report tables for one or more simulations}
\seeentry{cge_dynamic}{Solve a recursive dynamic run}
\seeentry{cge_gams}{Write the model as a GAMS program, and solve it in GAMS}
\seechapter{chap:microsim}{The link to household surveys}
\end{alsosee}
